\documentclass[10pt]{article}

\usepackage[utf8]{inputenc}
\usepackage[T1]{fontenc}
\usepackage{amsmath,amssymb,amsthm}
\usepackage[margin=0.85in]{geometry}
\usepackage{microtype}
\usepackage{hyperref}

\hypersetup{hidelinks}

\newtheorem{theorem}{Theorem}

\newcommand{\udens}{\overline{d}}
\newcommand{\ldens}{\underline{d}}

\title{\normalsize Note\\[1ex]
  \LARGE On permutations avoiding arithmetic progressions}
\author{Timothy D. LeSaulnier, Sujith Vijay\thanks{Corresponding author. E-mail address: \href{mailto:sujith@math.uiuc.edu}{sujith@math.uiuc.edu} (S. Vijay).}\\
  \small Department of Mathematics, University of Illinois at Urbana-Champaign,\\
  \small Urbana, IL 61801, USA}
\date{}

\begin{document}

\maketitle

\noindent
\textbf{Article history:}\par
\noindent
Received 2 June 2010\\
Received in revised form 14 August 2010\\
Accepted 11 October 2010\\
Available online 6 November 2010

\medskip
\noindent
\textbf{Keywords:}\par
\noindent
Arithmetic progressions\\
Permutations

\begin{abstract}
We improve the lower bound on the number of permutations of
$\{1,2,\ldots,n\}$ in which no 3-term arithmetic progression occurs as a
subsequence, and derive lower bounds on the upper and lower densities of
subsets of the positive integers that can be permuted to avoid 3-term and
4-term APs. We also show that any permutation of the positive integers must
contain a 3-term AP with odd common difference as a subsequence, and construct
a permutation of the positive integers that does not contain any 4-term AP
with odd common difference.
\end{abstract}

\begin{flushright}
\textcopyright{} 2010 Elsevier B.V. All rights reserved.
\end{flushright}

Let $S$ be a subset of the positive integers, and let $\sigma$ be a permutation
of $S$. We say that $\sigma$ is a $k$-avoiding permutation of $S$ if $\sigma$
does not contain any $k$-term AP as a subsequence. Similarly, the set $S$ is
said to be $k$-avoidable if there exists a $k$-avoiding permutation of $S$.

Let $M(n)$ denote the number of 3-avoiding permutations of
$\{1,2,\ldots,n\}$. For example, $M(4)=10$, corresponding to the permutations
$(1,3,4,2)$, $(1,3,2,4)$, $(2,1,4,3)$, $(2,4,1,3)$, $(4,2,1,3)$ and their
reversals. In 1977, Davis et al.\ [1] established the following bounds on
$M(n)$:
\[
  2^{n-1}
  \leq M(n)
  \leq \left\lfloor\frac{n+1}{2}\right\rfloor!
         \left\lceil\frac{n+1}{2}\right\rceil!.
\]
These bounds were recently improved by Sharma [3], who showed that
\[
  M(n) \leq \frac{2.7^n}{21}
  \qquad\text{for } n\geq 11,
\]
and that
\[
  \lim_{n\to\infty}\frac{M(n)}{2^n n^k}=\infty
  \qquad\text{for any fixed } k.
\]
In [3] the question whether
$\lim_{n\to\infty} M(n+1)/M(n)=2$ was attributed to the Editor of the Problem
Section of the \emph{American Mathematical Monthly} (where the function
$M(n)$ made its earliest known appearance, in 1975), and was mentioned as
still open. We begin with an observation that settles this question in the
negative. Indeed, we establish the following stronger lower bound on $M(n)$.

\begin{theorem}
$M(n)\geq (1/2)c^n$ for $n\geq 8$, where
$c=(2132)^{1/10}=2.152\ldots$.
\end{theorem}

\begin{proof}
The following inequalities were proved in [1] to show that
$M(n)\geq 2^{n-1}$:
\[
  M(2n)\geq 2[M(n)]^2;
  \qquad
  M(2n+1)\geq 2M(n)M(n+1).
\]
These recurrences follow from the observation that if $\sigma_1$ and
$\sigma_2$ are 3-avoiding permutations of $\{2,4,\ldots,2n\}$ and
$\{1,3,\ldots,2n-1\}$, respectively, concatenating them in either order
yields 3-avoiding permutations $\sigma_1\sigma_2$ and
$\sigma_2\sigma_1$ of $\{1,2,\ldots,2n\}$, since the first and third terms of
any arithmetic progression have the same parity. Note that these recurrences
imply the stronger lower bound $M(n)\geq (1/2)c^n$ for $n\geq 8$, where
$c=(2M(10))^{1/10}=2.152\ldots$. Since
\[
\begin{gathered}
  M(8)=282,\quad M(9)=496,\quad M(10)=1066,\quad M(11)=2460,\\
  M(12)=6128,\quad M(13)=12{,}840,\quad M(14)=29{,}380,\quad
  M(15)=74{,}904
\end{gathered}
\]
(see [1]), the inequality holds for $8\leq n\leq 15$. We can now use
induction on $k$ to show that it also holds for
$2^k\leq n<2^{k+1}$, $k\geq 4$.
\end{proof}

We now look at permutations of infinite subsets of the positive integers.
Davis et al.\ [1] observed that any permutation of the positive integers
contains a 3-term AP as a subsequence. (Let $a_1$ be the first term, and let
$k$ be the least integer such that $a_k>a_1$. Then $2a_k-a_1$ occurs to the
right of both $a_1$ and $a_k$.) They also constructed a 5-avoiding permutation
of the positive integers. The 4-avoidability of the positive integers remains
a fascinating open problem. However, if we restrict our attention to
arithmetic progressions with odd common difference, the question can be
answered.

\begin{theorem}
Any permutation of the positive integers must contain a 3-term AP with odd
common difference as a subsequence. Furthermore, there exists a permutation
of the positive integers in which no 4-term AP with odd common difference
occurs as a subsequence.
\end{theorem}

\begin{proof}
We first show that any permutation
$\sigma=(t_1,t_2,\ldots,t_{11})$ of $\{1,2,\ldots,11\}$ with $t_1=2$ and
$t_2=1$ must contain a 3-term AP with odd common difference as a subsequence.
Indeed, 4 must appear in $\sigma$ before 3 to avoid the 3-term AP $(2,3,4)$;
similar considerations force 4 before 5, 7 before 4, 6 before 5, 6 before 7,
11 before 6, 8 before 11, 8 before 9 and 7 before 10. Thus, both 8 and 11
appear in $\sigma$ before either 9 or 10 appears. Now we have the 3-term AP
$(8,9,10)$ if 9 appears before 10 in $\sigma$ and the 3-term AP $(11,10,9)$
otherwise. This proves our claim.

Let $a_1,a_2,\ldots$ be a permutation of the positive integers. We may assume
without loss of generality (subtracting $a_1-1$ from each term and ignoring
non-positive terms) that $a_1=1$. Let $k$ be the least index such that $a_k$
is even, and let
\[
  a_j=\max(a_1,a_2,\ldots,a_{k-1}).
\]
If $a_j<2a_k-1$, then we have $(1,a_k,2a_k-1)$ as a subsequence. If
$a_j\geq 2a_k-1>a_k$, then let $d=a_j-a_k$. Note that $d$ is odd. Since
$a_j$ occurs before $a_k=a_j-d$ and they both occur before any of
$a_j+d,a_j+2d,\ldots,a_j+9d$, it follows from the above claim (via shifting
and scaling) that the permutation contains a 3-term AP with odd common
difference.

Now we exhibit a permutation of the positive integers that contains no
4-term AP with odd common difference as a subsequence. For $i\geq 1$, let
$\sigma_i$ be a 3-avoiding permutation of the following set of
$2^{2i-1}$ consecutive even numbers:
\[
  \left\{
    \frac{4^i+2}{3},\frac{4^i+8}{3},\ldots,
    \frac{4^{i+1}-4}{3}
  \right\}.
\]
Similarly, let $\pi_i$ be a 3-avoiding permutation of the following set of
$4^{i-1}$ consecutive odd numbers:
\[
  \left\{
    \frac{4^i+2}{6},\frac{4^i+14}{6},\ldots,
    \frac{4^{i+1}-10}{6}
  \right\}.
\]
Observe that the concatenated sequence
$\sigma_1\pi_1\sigma_2\pi_2\sigma_3\pi_3\cdots$ is a permutation of the
positive integers. By virtue of our construction, if an odd number $x$
occurs in this sequence before an even number $y$, then $2x-y<0$. It follows
that no 4-term AP with odd common difference occurs as a subsequence.
\end{proof}

Given a subset $S$ of the positive integers, let $\udens(S)$ and $\ldens(S)$
denote, respectively, the upper and lower densities of $S$. In other words,
\[
  \udens(S)=\limsup_{n\to\infty}\frac{S(n)}{n}
  \qquad\text{and}\qquad
  \ldens(S)=\liminf_{n\to\infty}\frac{S(n)}{n},
  \qquad
  S(n)=\lvert S\cap[1,n]\rvert.
\]
Define, for $k\geq 3$,
\[
\begin{aligned}
  \alpha(k)&=\sup_S\{\udens(S):S\text{ is }k\text{-avoidable}\},\\
  \beta(k)&=\sup_S\{\ldens(S):S\text{ is }k\text{-avoidable}\}.
\end{aligned}
\]
Since the set of positive integers is 5-avoidable,
$\alpha(k)=\beta(k)=1$ for $k\geq 5$. Bounds for $\alpha(3)$ and $\beta(3)$
were sought in [1]. We show the following.

\begin{theorem}
$\alpha(4)=1$, $\alpha(3)\geq 1/2$, $\beta(4)\geq 1/3$ and
$\beta(3)\geq 1/4$.
\end{theorem}

\begin{proof}
For integers $a\geq 2$ and $i\geq 0$, define
\[
  S_i^{(a)}=\{a^{2i},a^{2i}+1,\ldots,a^{2i+1}\},
\]
and let $\sigma_i^a$ be a 3-avoiding permutation of $S_i^{(a)}$. Define
\[
  S^{(a)}=\bigcup_{i\geq 0}S_i^{(a)}.
\]
We claim that $S^{(a)}$ is 4-avoidable. Clearly, the concatenated sequence
$\sigma_0^a\sigma_1^a\cdots$ does not contain a decreasing 3-term AP. Suppose
it contains an increasing 4-term AP $x_1,x_2,x_3,x_4$. Since $x_2,x_3$ and
$x_4$ cannot all belong to the same set $S_i^{(a)}$, we must have
$x_4\geq 2x_3$ or $x_3\geq 2x_2$. But then $x_2\leq 0$ or $x_1\leq 0$,
yielding a contradiction. Note that
\[
\begin{aligned}
  \udens\bigl(S^{(a)}\bigr)
    &=\frac{a-1}{a}\sum_{i=0}^{\infty}a^{-2i}
      =\frac{a}{a+1},\\
  \ldens\bigl(S^{(a)}\bigr)
    &=\frac{a-1}{a^2}\sum_{i=0}^{\infty}a^{-2i}
      =\frac{1}{a+1}.
\end{aligned}
\]
Since $a$ can be arbitrarily large, it follows that $\alpha(4)=1$. Taking
$a=2$, we get $\beta(4)\geq 1/3$.

Let $p_0=1$, $q_0=2$, and, for $k\geq 1$, define
$p_k=2q_{k-1}$ and $q_k=3q_{k-1}-1$. Let $\tau_k$ be a 3-avoiding
permutation of
\[
  T_k=\{p_k,p_k+1,\ldots,q_k\},
\]
and let
\[
  T=\bigcup_{k\geq 0}T_k.
\]
Since $p_k=3^k+1=2q_{k-1}$ for $k\geq 1$, it follows that
$\udens(T)=1/2$ and $\ldens(T)=1/4$. We claim that the concatenated sequence
$\tau_0\tau_1\cdots$ is a 3-avoiding permutation of $T$. Indeed, if the
(increasing) 3-term AP $x_1,x_2,x_3$ occurs as a subsequence, with $x_2$ and
$x_3$ belonging to different sets $T_k$ and $T_\ell$, then $x_3\geq 2x_2$,
so $x_1\leq 0$, yielding a contradiction. If $x_2$ and $x_3$ belong to the
same set $T_k$, then $x_1\in T_\ell$ with $\ell<k$. But
\[
  x_3-x_2<q_{k-1}\leq x_2-x_1,
\]
contradicting our assumption that $x_1,x_2,x_3$ is a 3-term AP. Therefore,
$T$ is 3-avoidable. Thus $\alpha(3)\geq 1/2$ and $\beta(3)\geq 1/4$.
\end{proof}

Erd\H{o}s and Graham [2] (see also [1]) asked if the positive integers can be
partitioned into two 3-avoidable sets. Clearly, the answer is negative if
$\alpha(3)+\beta(3)<1$. We believe this to be the case, and conjecture that
the lower bounds in the above theorem are optimal, i.e., $\alpha(3)=1/2$ and
$\beta(3)=1/4$. However, we have not even been able to show that
$\beta(3)<1$.

\section*{Acknowledgements}

We thank the referees for several valuable suggestions, which improved the
overall readability of the paper.

\section*{References}

\begin{enumerate}
\item J.A. Davis, R.C. Entringer, R.L. Graham, G.J. Simmons,
  On permutations containing no long arithmetic progressions,
  \emph{Acta Arithmetica} 34 (1977), 81--90.

\item P. Erd\H{o}s, R.L. Graham,
  \emph{Old and New Problems and Results in Combinatorial Number Theory},
  Monographie No.~28 de L'Enseignement Math\'ematique, Geneva, 1980.

\item A. Sharma,
  Enumerating permutations that avoid three term arithmetic progressions,
  \emph{Electronic Journal of Combinatorics} 16 (2009), \#R63.
\end{enumerate}

\end{document}
